\subsection*{Chapter 5: Titrimetric Methods of Analysis}

\subsubsection*{1. Calculation of Results}
% Source: Chapter 5.1, Slides 21-23

\textbf{Normality Calculation:}
\begin{equation}
C_{N(X)} = \frac{C_C V_C}{V_X}
\end{equation}

\textbf{Mass Concentration:}
\begin{equation}
C(g/l) = C_{N(X)} \times E_X
\end{equation}

\textbf{Percentage in Solid Sample:}
\begin{equation}
\%X = C_C \times V_C \times 10^{-3} \times \frac{100}{m} \times E_X
\end{equation}

\textbf{Titer ($T$):}
\begin{equation}
T_{C/X} = C_C \times 10^{-3} \times E_X
\end{equation}
\textit{$T_{C/X}$: grams of analyte X reacting with 1 mL of titrant C.}

\end{multicols}

\begin{table}[htbp]
    \centering
    \caption{Overview of Titration Methods}
    \small % Reduce font size slightly to ensure perfect fit
    \renewcommand{\arraystretch}{1.3} % Add breathing room between rows
    
    % Total width approx 18.5cm. 
    % Columns: Type, Method, Titrant, Analyte, Indicator, Color Change, Notes
    \begin{tabularx}{\textwidth}{@{} L{1.7cm} L{2.2cm} L{1.4cm} L{1.8cm} L{2.4cm} L{2.8cm} Y @{}}
        \toprule
        \textbf{Type} & \textbf{Method} & \textbf{Titrant} & \textbf{Analyte} & \textbf{Indicator} & \textbf{Color Change} & \textbf{Conditions / Notes} \\
        \midrule
        
        % --- REDOX ---
        \textbf{Redox} 
        & Permanganate & \ce{KMnO4} & \ce{H2C2O4}, \ce{Fe^{2+}} 
        & Self-indicating & Purple (\ce{MnO4^-}) $\to$ Colorless (\ce{Mn^{2+}}) & Acidic medium. Heat to $\approx 70^\circ$C often required. \\ 
        
        & Dichromate & \ce{K2Cr2O7} & \ce{Fe^{2+}} 
        & Diphenylamine & Colorless $\to$ Purple & Acidic medium. \\ 
        
        & Iodimetry & \ce{I2} & Reducers & Starch & Colorless $\to$ Blue & Direct titration. Blue appears at endpoint. \\ 
        
        & Iodometry & \ce{Na2S2O3} & Oxidizers (\ce{Cu^{2+}}) 
        & Starch & Blue $\to$ Colorless & Indirect. \ce{I2} released is titrated; blue disappears. \\ 
        \midrule
        
        % --- ACID-BASE ---
        \textbf{Acid-Base} 
        & Strong/Strong & \ce{NaOH}/\ce{HCl} & \ce{HCl}/\ce{NaOH} 
        & Phenolphthalein & Colorless $\to$ Pink & Equivalence point at pH 7. \\ 
        
        & Weak Acid & \ce{NaOH} & \ce{CH3COOH} 
        & Phenolphthalein & Colorless $\to$ Pink & Equivalence point at pH $>7$. \\ 
        
        & Polyprotic Acid & \ce{NaOH} & \ce{H3PO4} 
        & 1. Methyl red \newline 2. Phenolphthalein & 1. Red $\to$ Yellow \newline 2. Colorless $\to$ Pink & Stepwise titration. Separate endpoints. \\ 
        
        & Polyprotic Base & \ce{HCl} & \ce{CO3^{2-}} 
        & 1. Phenolphthalein \newline 2. Methyl orange & 1. Pink $\to$ Colorless \newline 2. Yellow $\to$ Orange & Step 1: \ce{CO3^{2-} -> HCO3^-}. \newline Step 2: \ce{HCO3^- -> H2CO3}. \\ 
        \midrule
        
        % --- PRECIPITATION ---
        \textbf{Precipitation} 
        & Mohr's Method & \ce{AgNO3} & Halogens (\ce{Cl^-}) 
        & Chromate & Yellow $\to$ Red ppt & Direct. pH 6.5--10. \ce{Ag2CrO4} forms. \\ 
        
        & Fajans' Method & \ce{AgNO3} & Halogens (\ce{Cl^-}) 
        & Adsorption Ind. & Yellow/Green $\to$ Pink & Direct. pH 6.5--10. Ind. adsorbs to precipitate. \\ 
        
        & Volhard's & \ce{KSCN} & \ce{Ag^+} (excess) 
        & \ce{Fe^{3+}} (Alum) & Colorless $\to$ Red & Back titration. Acidic (pH $<3$). \ce{Fe(SCN)^{2+}}. \\ 
        \midrule
        
        % --- COMPLEXATION ---
        \textbf{Complexation} 
        & Total Hardness & EDTA & \ce{Ca^{2+}} + \ce{Mg^{2+}} 
        & Erio Black T & Wine-red $\to$ Blue & Buffer pH 10. \\ 
        
        & Calcium Only & EDTA & \ce{Ca^{2+}} 
        & Fluorexone & Green fl. $\to$ Pink & pH 12.5. \ce{Mg^{2+}} precipitates as \ce{Mg(OH)2}. \\ 
        
        & Iron (Direct) & EDTA & \ce{Fe^{3+}} 
        & Sulfosalicylic Acid & Violet $\to$ Colorless/Yellow & pH 2.5. Heat gently. \\ 
        
        & Aluminum (Back) & \ce{Cu^{2+}}/\ce{Zn^{2+}} & \ce{Al^{3+}} 
        & PAN & Yellow $\to$ Violet & Back titration. Excess EDTA added at pH 5, heated. \\ 
        
        \bottomrule
    \end{tabularx}
\end{table}

\begin{multicols}{2}
\subsubsection*{2. Acid-Base Titrations}

\textbf{Polyprotic Acids (Ampholytes):}
\begin{equation}
pH_{eq1} = \frac{1}{2}(pK_{a1} + pK_{a2})
\end{equation}
\textit{(Valid for the first equivalence point of acids like $H_3PO_4$)}

\subsubsection*{3. Redox Titrations}
\textbf{Potential using Fraction Titrated ($F$):}
\begin{equation}
E_{sol} = E^0_X + \frac{0.059}{n_X} \log \frac{F}{1-F}
\end{equation}

\subsubsection*{4. Indicator Error}
% Source: Chapter 5.3

\textbf{General Indicator Error Equation:}
\begin{equation}
\Delta\% = \frac{[X]_{\text{equiv}}(10^{\Delta pX} - 10^{-\Delta pX})}{[X]_0} \times 100
\end{equation}
Where \(\Delta pX = |pX_f - pX_{\text{equiv}}|\).

\textbf{Complexation Titration Error:}
\begin{equation}
\Delta\% = \frac{10^{\Delta pX} - 10^{-\Delta pX}}{([X]_0 \cdot \beta_{XC})^{1/2}} \times 100
\end{equation}

\textbf{Precipitation Titration Error:}
\begin{equation}
\Delta\% = \frac{10^{\Delta pX} - 10^{-\Delta pX}}{([X]_0^2 / T_{CX})^{1/2}} \times 100
\end{equation}

\textbf{Strong Acid - Strong Base Error (Approximate):}
% Source: Chapter 5.3, Slide 31
\begin{equation}
\Delta\% = \frac{10^{-pT}(C_X + C_C)}{C_X C_C} \times 100
\end{equation}
\textit{(Where $pT$ is the pH difference from neutral)}

\textbf{Redox Titration Error (Approximate) - 3 Step Method:}

\textbf{Step 1: Determine Curve Direction}
\begin{itemize}
    \item \textbf{Upward Curve:} Analyte (X) is Reducing Agent ($E^0_X < E^0_C$).
    \item \textbf{Downward Curve:} Analyte (X) is Oxidizing Agent ($E^0_X > E^0_C$).
\end{itemize}

\textbf{Step 2: Compare $E_f$ vs $E_{equiv}$}
\begin{itemize}
    \item Calculate $E_f$ (Indicator limit):
    \begin{itemize}
        \item \textbf{Upward Curve:} $E_f = E^0_{ind} + \frac{0.059}{n_{ind}}$
        \item \textbf{Downward Curve:} $E_f = E^0_{ind} - \frac{0.059}{n_{ind}}$
    \end{itemize}
    \item Calculate $E_{equiv}$.
\end{itemize}

\textbf{Step 3: Select Formula}

\noindent
\small
\begin{tabularx}{\linewidth}{@{} L{1.2cm} Y Y @{}}
\toprule
 & \textbf{A. Upward} \newline ($E^0_X < E^0_C$) & \textbf{B. Downward} \newline ($E^0_X > E^0_C$) \\
\midrule
\textbf{Before Eq.} & \textbf{($E_f < E_{equiv}$)} \newline $\Delta\% = 10^{\frac{-n_X(E_f - E_X^0)}{0.059}} \times 100$ & \textbf{($E_f > E_{equiv}$)} \newline $\Delta\% = 10^{\frac{n_X(E_f - E_X^0)}{0.059}} \times 100$ \\
\addlinespace[0.5em]
\midrule
\addlinespace[0.5em]
\textbf{After Eq.} & \textbf{($E_f > E_{equiv}$)} \newline $\Delta\% = 10^{\frac{n_C(E_f - E_C^0)}{0.059}} \times 100$ & \textbf{($E_f < E_{equiv}$)} \newline $\Delta\% = 10^{\frac{-n_C(E_f - E_C^0)}{0.059}} \times 100$ \\
\bottomrule
\end{tabularx}
\normalsize

